\begin{code}[hide]
{-# OPTIONS --cubical-compatible --safe #-}

open import Data.Fin using (toℕ)
open import Data.Nat using (ℕ ; suc)
open import Data.Nat.Primality using (Prime)
open import Data.Product using (_,_ ; ∃)
open import Relation.Binary.PropositionalEquality using (_≡_)

open import Notations
open import ForStdlib.Data.Fin.Mod
open import ForStdlib.Data.Fin.Mod.Prime.Fermat

module vqupit.sections.results
  (p-3 : ℕ)
  (let p-2 = ₁₊ p-3)
  (p-prime : Prime (suc (₁₊ p-2)))
  (let open PrimeModulus' p-2 p-prime)
  (g*@(g , g≠0) : ℤ* ₚ)
  (g-gen : ∀ ((x , _) : ℤ* ₚ) → ∃ λ (k : ℤ ₚ-₁) → x ≡ g ^′ toℕ k)
  where

open import Algebra.Bundles using (Group)
open import Word.Base using (_•_)
open import Presentation.Construct.Base using ([_]ₗ ; [_]ᵣ)
open import Presentation.Definitions using (_IsPresentationOf_)
import Presentation.Base as PB
import Examples.Groups.Qupit-Clifford.Syntactics
  p-3 p-prime g* g-gen as Syn
import Examples.Groups.Qupit-Clifford.Semantics p-3 p-prime as Model
import Examples.Groups.Qupit-Clifford.Structural
  p-3 p-prime g* g-gen as Structural
import Examples.Groups.Qupit-Clifford.Certificate
  p-3 p-prime g* g-gen as Client
import Examples.Groups.Qupit-Clifford.SemFE
  p-3 p-prime g* g-gen as FE
import Examples.Groups.Cyclic.Scalars as CyclicScalars
import Examples.Groups.Clifford+MinusOne.Qupit.Figure1.Syntactics
  p-3 p-prime g* g-gen as F
import Examples.Groups.Clifford+MinusOne.Qupit.Figure1.Structural
  p-3 p-prime g* g-gen as Signed

module Sc = CyclicScalars p-2
open Syn using (_Exact,_===_)
open F using (_F,_===_)
open Signed using (StructuralClifford±)
\end{code}
\subsection{Structural semantics and a certificate client}\label{sec:structural}

Write $J_p(n)$ for the structural Weyl extension. Its elements are
$(e,P,S)\in\Zp\times\Zp^{2n}\times\Sp$, with multiplication
\[
(e,P,S)(f,Q,T)=
\bigl(e+f+\tfrac12\mathrm{sform}(P,SQ),\ P+SQ,\ ST\bigr).
\]
This definition contains no circuit congruence. To connect it to
\aid{Exact-group}, define $\chi(a,u)=(\mathrm{read}(a),\sem{u}_{\rm proj})$,
where $a$ is a scalar word and $u$ a projective gate word.

\begin{theorem}[Structural scalar presentation]\label{thm:structural}
For every width $n$, $\chi$ is a group isomorphism from the presented
scalar extension to $J_p(n)$. Composing it with the exact presentation
gives a presentation of $J_p(n)$ with the specified interpretation.
\end{theorem}
\begin{proof}
The scalar presentation reads $\omega^k$ as $k\bmod p$, and the
projective presentation reads $u$ as its Pauli and symplectic components.
The generated cocycle is congruent to the scalar-word embedding of the
pulled-back Weyl cocycle (\aid{SemGeneratorData.agree}).
Reading this equation in $\Zp$ proves that $\chi$ preserves the twisted
multiplication. Injectivity and surjectivity follow componentwise from
the two factor presentations. This also fixes the interpretation:
scalar words retain their exponents, and gate words retain their
projective interpretation and accumulated Weyl phase.
\end{proof}

\begingroup\renewcommand{\AgdaCodeStyle}{\footnotesize}
\begin{code}
exact-structural-presentation : ∀ n →
  (n Exact,_===_) IsPresentationOf (Model.Clifford-group n)
exact-structural-presentation = Structural.Bridge.structural-presentation
\end{code}
\endgroup

\begin{theorem}[Original scalar-sensitive rules]\label{thm:figure1}
For every width $n$, the original sixteen rules over $-\omega,H,S,CZ$
present $J_p(n)\times\mathbb Z_2$.
\end{theorem}
The map is $\chi$ on the exact component and the identity on the sign.
Composing with the checked Figure~1 translation gives:
\begingroup\renewcommand{\AgdaCodeStyle}{\footnotesize}
\begin{code}
figure1-structural-presentation : ∀ n →
  (n F,_===_) IsPresentationOf (StructuralClifford± n)
figure1-structural-presentation = Signed.At.presentation
\end{code}
\endgroup
These excerpts are aliases of the actual implementations, with no
postulated abbreviations; their literate module is also fully typechecked.
The parameters are the odd prime and primitive-root witness of
\S\ref{sec:symplectic}. The result strengthens the semantic endpoint,
but does not assert that the named generators denote the intended
complex matrices.

\paragraph{Repairing a projectively valid replacement.}
Suppose a client proposes to replace $u$ by $v$ and proves
$\sem u_{\rm proj}=\sem v_{\rm proj}$. This premise includes the Pauli
component: agreement of symplectic matrices alone is insufficient.
The Weyl evaluator computes the scalar phase by
\[
\Phi(\varepsilon)=\Phi([a])=0,\qquad
\Phi(uv)=\Phi(u)+\Phi(v)+c(\sem u_{\rm proj},\sem v_{\rm proj}).
\]
Define $t=\Phi(u)-\Phi(v)$ in $\Zp$. The client returns the exact
derivation
\[
 [u]_{\rm gate}\ \approx_{\rm exact}\
 [\omega^t]_{\rm scalar}\bullet[v]_{\rm gate}.
\]
Its implementation reads both sides in $J_p(n)$, proves equality by
finite-field arithmetic and the supplied projective equality, and
applies the presentation's injectivity. It performs no search through
the circuit equations.

\begingroup\renewcommand{\AgdaCodeStyle}{\footnotesize}
\begin{code}
phase-repair : ∀ n (u v : Syn.Circuit n) →
  Group._≈_ (Model.Pauli⋊Sp-group n)
    (FE.Weyl.⟦_⟧ n u) (FE.Weyl.⟦_⟧ n v) →
  PB._≈_ (n Exact,_===_) [ u ]ᵣ
    ([ Sc.emb (Client.At.correction n u v) ]ₗ • [ v ]ᵣ)
phase-repair = Client.At.repair
\end{code}
\endgroup

\paragraph{A checked two-qutrit replacement.}\label{sec:client-example}
At $p=3$, let $H_i,S_i$ act on wire $i$, with wire zero at the bottom,
and write $C=CZ_{0,1}$. Products use the circuit syntax's order. Consider
\[
\begin{aligned}
u &= S_1 C H_1 H_0 H_1 C S_1 S_0 H_0 C S_0 H_0,\\
v &= S_0 S_0 S_1 H_1 H_1 C.
\end{aligned}
\]
The module \texttt{Qupit-Clifford.CertificateExample} supplies the
complete semantic premise and calls \aid{repair}. A compact evaluator
computes the Pauli vector, symplectic action, and Weyl phase; a structural
induction proves agreement with the presentation's interpretation.
Agda checks equality of the Pauli vectors, checks the actions on all
$3^4=81$ phase-space vectors, and computes
$\Phi(u)=1$, $\Phi(v)=0$. The result is the exact certificate
$[u]_{\rm gate}\approx_{\rm exact}
[\omega]_{\rm scalar}\bullet[v]_{\rm gate}$,
which \aid{replace} lifts into arbitrary surrounding words.

The gate counts are also checked: twelve $H,S,CZ$ gates become six,
and three entangling gates become one. The replacement retains one
explicit scalar gate; the counts exclude that scalar. Deleting it
would lose the exact phase. An external bounded search found this
candidate; \path{notes/find-certificate-example.py} reproduces the
search. Its output is untrusted: the Agda module separately verifies
the equality and constructs the derivation. This example exercises
semantic checking, phase repair, and contextual replacement together;
it is not a benchmark or an optimality result.

For the original gate language, \aid{Figure1.Structural.certify} takes
equality in $J_p(n)\times\mathbb Z_2$ and returns a derivation in its
sixteen rules. Its \aid{replace} operation lifts that derivation into
any surrounding words $l,r$. These interfaces demonstrate certificate
construction, not equivalence discovery: the caller must supply the
semantic equality proof. There is no packaged full-circuit decision
procedure, optimizer, or claim about certificate size.
